\documentclass[12pt,a4paper]{article}

\usepackage[utf8]{inputenc}
\usepackage[T1]{fontenc}
\usepackage[english]{babel}
\usepackage[margin=2.5cm]{geometry}
\usepackage{mathpazo}

\title{\textbf{Performance Optimization of the AI4Land Training Pipeline on MareNostrum~5}}
\author{
  Ferran Mirabent Rubinat\\[1em]
  \small Master's Thesis in Modelling for Science and Engineering\\
  \small Universitat Aut\`onoma de Barcelona\\[1.5em]
  \small \textbf{Tutors:}\\
  \small Oriol Tint\'o Prims \\
  \small Barcelona Supercomputing Center (BSC)\\\and
  \small S\'ilvia Cuadrado Gavil\'an\\[0.5em]
  \small Universitat Aut\`onoma de Barcelona (UAB)
}
\date{March 2026}

\begin{document}

\maketitle
\thispagestyle{empty}

\begin{abstract}
Accurate modelling of land use and land cover (LULC) change is essential for projecting carbon
emissions and understanding climate feedbacks. The AI4Land project, developed at the Barcelona
Supercomputing Center within the framework of the Horizon Europe projects CONCERTO and TerraDT,
employs a UNet-based deep learning model to downscale LULC data from the Land-Use Harmonization
dataset (LUH2) at approximately 30\,km resolution to 1\,km resolution, covering the period
1850--2100 including future projections based on Shared Socioeconomic Pathways (SSPs). The
pipeline is also intended to be adopted by the ELLIOT project for training geospatial foundation
models. Training this model on multi-channel geospatial datasets is computationally demanding
and requires efficient use of high-performance computing resources.

This thesis focuses on identifying and resolving performance bottlenecks in the AI4Land training
pipeline running on MareNostrum~5, equipped with NVIDIA H100 GPUs. Preliminary profiling has
revealed the data loading stage as the primary bottleneck, with significant variance in epoch
duration traceable to the current use of multiple Zarr stores and suboptimal data ingestion
patterns. The work will systematically evaluate optimizations with a particular emphasis on the
data loading pipeline, while also exploring computational improvements at the model and hardware
level. Each optimization will be assessed through profiling with NVIDIA Nsight Systems to quantify
its impact on throughput and training time.

The expected outcome is a measurably faster and more efficient training pipeline that maintains
model accuracy, enabling faster iteration cycles for the AI4Land project and contributing
practical guidelines for training deep learning models on large-scale HPC infrastructure.
\end{abstract}

\end{document}
